% partADC_3a.tex --- Phan 3a: Adaptive Density Control -- E_photo va lien he voi Pruning
% 6 frame, khong documentclass/usepackage/begin{document} (duoc include tu main.tex)

% --- 3_1_ephoto-3view.png ---
\begin{frame}{Bước (6): tổng hợp sai số ảnh $E_{\text{photo}}$ qua nhiều view}
\begin{columns}[T]
\begin{column}{0.48\linewidth}
\footnotesize
\begin{itemize}
  \item Mỗi view $v$ cho ra một cặp $(L1^{(v)}, \text{SSIM}^{(v)})$ giữa ảnh render và ảnh GT.
  \item Công thức tổng hợp (code: \texttt{e\_photo}, $\lambda = 0.2$):
  \[
    E_{\text{photo}}^{(v)} = (1-\lambda)\, L1^{(v)} + \lambda\,\bigl(1 - \text{SSIM}^{(v)}\bigr)
  \]
  \item Đây chính là dạng \emph{trung bình có trọng số} giữa hai thành phần L1 và $1-\text{SSIM}$, không phải max qua view.
  \item Ví dụ 3.md (ảnh 4x4, 3 view): $L1 = (0.11125,\, 0.06,\, 0.090625)$, $\text{SSIM} = (0.72,\, 0.88,\, 0.80)$.
  \item Kết quả: $E_{\text{photo}} = (0.145,\; 0.072,\; 0.1125)$ -- view 1 (SSIM thấp nhất) đóng góp lỗi lớn nhất.
\end{itemize}
\end{column}
\begin{column}{0.5\linewidth}
\centering
\includegraphics[width=0.95\linewidth]{3_1_ephoto-3view.png}
\captionof{figure}{\footnotesize Phân rã $E_{\text{photo}}^{(v)} = 0.8\,L1 + 0.2(1-\text{SSIM})$ cho 3 view, cột xanh lá = giá trị dùng làm trọng số ở bước (7).}
\end{column}
\end{columns}
\end{frame}

% --- 3_2_l1-vs-ssim.png ---
\begin{frame}{Vì sao cần SSIM: L1 không phân biệt được ``mờ'' và ``nhiễu''}
\begin{columns}[T]
\begin{column}{0.52\textwidth}
\scriptsize
\begin{itemize}
  \item Thí nghiệm: GT tổng hợp, Render A (blur Gaussian $\sigma{=}1.6$), Render B (nhiễu hiệu chỉnh để \textbf{cùng} $L1$ với A).
  \item $L1^{A}{\approx}L1^{B}$ nhưng cấu trúc khác $\Rightarrow \text{SSIM}^{A}{\neq}\text{SSIM}^{B}$ (blur phá cạnh giữ tương quan patch, noise phá cả hai).
  \item $\Rightarrow E_{\text{photo}}^{A}{\neq}E_{\text{photo}}^{B}$ dù $L1$ gần bằng nhau — chỉ dùng $L1$ sẽ coi hai render "lỗi như nhau".
  \item[$\Rightarrow$] $E_{\text{photo}}$ cần cả L1 (cường độ) và $1{-}$SSIM (cấu trúc, cửa sổ Gauss $11{\times}11$).
\end{itemize}
\end{column}
\begin{column}{0.48\textwidth}
\centering
\includegraphics[width=0.95\linewidth]{3_2_l1-vs-ssim.png}
\captionof{figure}{\footnotesize Cùng $L1$, SSIM và $E_{\text{photo}}$ khác nhau giữa render mờ và nhiễu.}
\end{column}
\end{columns}
\end{frame}

% --- 3_3_ephoto-theo-lambda.png ---
\begin{frame}{$E_{\text{photo}}$ theo trọng số trộn $\lambda$}
\begin{columns}[T]
\begin{column}{0.48\linewidth}
\footnotesize
\begin{itemize}
  \item $E_{\text{photo}}^{(v)}(\lambda) = (1-\lambda)\,L1^{(v)} + \lambda\,(1-\text{SSIM}^{(v)})$ là hàm \emph{tuyến tính} theo $\lambda \in [0,1]$.
  \item $\lambda = 0$: chỉ dùng $L1$ thuần; $\lambda = 1$: chỉ dùng $1-\text{SSIM}$ (giống vai trò $\lambda_{\text{dssim}}$ trong loss huấn luyện).
  \item Giá trị mặc định trong code: $\lambda = 0.2$ -- ưu tiên L1 nhưng vẫn giữ một phần tín hiệu cấu trúc.
  \item Các đường $E_{\text{photo}}^{(v)}(\lambda)$ của từng view có thể \textbf{cắt nhau}: view có $L1$ thấp nhưng SSIM thấp có thể vượt lên khi $\lambda$ tăng.
  \item Tỉ số $E_{\max}/E_{\min}$ qua các view thay đổi rõ rệt giữa $\lambda=0$ và $\lambda=1$, cho thấy lựa chọn $\lambda$ ảnh hưởng trực tiếp đến việc view nào bị coi là ``tệ nhất''.
\end{itemize}
\end{column}
\begin{column}{0.5\linewidth}
\centering
\includegraphics[width=0.95\linewidth]{3_3_ephoto-theo-lambda.png}
\captionof{figure}{\footnotesize $E_{\text{photo}}^{(v)}(\lambda)$ cho 3 view (ví dụ 3.md và cảnh đồ chơi); đường đứt nét đánh dấu $\lambda=0.2$ dùng trong code.}
\end{column}
\end{columns}
\end{frame}

% --- 3_4_heatmap-dong-gop.png ---
\begin{frame}{Đóng góp của từng Gaussian vào điểm số tổng}
\footnotesize
\begin{itemize}
  \item Với mỗi Gaussian $i$ và view $v$, \texttt{counts}$_i^{(v)}$ (bước 4) đếm số pixel lỗi mà $i$ tham gia dựng nên trong view đó.
  \item Đóng góp có trọng số vào điểm thô: $\text{contrib}_i^{(v)} = \text{counts}_i^{(v)} \cdot E_{\text{photo}}^{(v)}$.
  \item Điểm thô tích luỹ qua $V$ view: $\text{raw}_i = \sum_{v=1}^{V} \text{counts}_i^{(v)} \, E_{\text{photo}}^{(v)}$.
  \item Ví dụ 3.md: $G3$ có \texttt{counts} $=(6,6,6)$ ở cả 3 view $\Rightarrow \text{raw}_{G3} = 6(0.145+0.072+0.1125) \approx 1.98$ -- cao nhất trong 5 Gaussian, trong khi $G1$ với \texttt{counts}$=(0,0,0)$ có $\text{raw}=0$.
  \item Heatmap minh hoạ trực quan: chỉ một vài Gaussian (hàng có màu đậm) chiếm phần lớn tổng sai số ảnh, phần còn lại đóng góp gần như bằng 0.
\end{itemize}
\centering
\includegraphics[width=0.85\linewidth]{3_4_heatmap-dong-gop.png}
\captionof{figure}{\footnotesize Trái: \texttt{counts}$_i^{(v)}$; giữa: đóng góp có trọng số \texttt{counts}$\cdot E_{\text{photo}}$; phải: điểm thô và Pruning sau minmax qua 5 Gaussian.}
\end{frame}

% --- 3_5_scatter-imp-pruning.png ---
\begin{frame}{Importance vs.\ Pruning: ranh giới không phải ngưỡng cứng}
\begin{columns}[T]
\begin{column}{0.48\linewidth}
\footnotesize
\begin{itemize}
  \item Importance: $\text{Importance}_i = \left\lfloor \dfrac{1}{V}\sum_{v=1}^{V} \text{counts}_i^{(v)} \right\rfloor$ -- số nguyên, dùng để quyết định \emph{densify} (ngưỡng $> 5$).
  \item Pruning: $\text{Pruning}_i = \text{minmax}_i(\text{raw}_i)$, chuẩn hoá $\text{raw}$ về $[0,1]$ qua toàn bộ $N$ Gaussian -- dùng cho \emph{final prune} (ngưỡng $>0.9$).
  \item Trên tập mô phỏng 300 Gaussian, $V{=}10$: hai đại lượng tương quan hạng dương (hệ số Spearman $\rho$ hiển thị trên hình) nhưng \textbf{không phải hàm một-một} của nhau.
  \item Với cùng một giá trị Importance, Pruning trải rộng trên một khoảng giá trị -- do Pruning còn phụ thuộc \emph{view nào} chứa lỗi (qua trọng số $E_{\text{photo}}^{(v)}$), còn Importance chỉ đếm số lượt xuất hiện.
  \item Kết luận: quyết định pruning có vùng chuyển tiếp mờ quanh ranh giới, không cắt cứng theo Importance.
\end{itemize}
\end{column}
\begin{column}{0.5\linewidth}
\centering
\includegraphics[width=0.95\linewidth]{3_5_scatter-imp-pruning.png}
\captionof{figure}{\footnotesize Trái: 5 Gaussian ví dụ 3.md; phải: 300 Gaussian mô phỏng, màu = tỉ lệ view thấy được Gaussian.}
\end{column}
\end{columns}
\end{frame}

% --- 3_6_pruning-vs-1-minus-imp.png ---
\begin{frame}{Pruning $\neq 1 - $ Importance chuẩn hoá}
\footnotesize
\begin{itemize}
  \item Bốn kịch bản cùng tổng \texttt{counts} nhưng phân bố khác nhau qua 3 view: A $=(6,6,6)$, B $=(18,0,0)$, C $=(0,18,0)$, D $=(17,0,0)$.
  \item A, B, C có cùng $\text{Importance} = \lfloor 18/3 \rfloor = 6$, nhưng điểm thô $\text{raw} = \sum_v \text{counts}^{(v)} E_{\text{photo}}^{(v)}$ khác nhau rõ rệt: dồn lỗi vào view có $E_{\text{photo}}$ cao (view 1, $E{=}0.145$) cho \text{raw} lớn hơn dồn vào view tốt nhất (view 2, $E{=}0.072$).
  \item D có tổng thấp hơn (17) nên Importance có thể giảm bậc trong khi raw vẫn gần B -- cho thấy việc lấy $\lfloor \cdot \rfloor$ làm mất thông tin mà Pruning (không làm tròn, có trọng số $E_{\text{photo}}$) vẫn giữ được.
  \item Trên 300 Gaussian mô phỏng: ứng với mỗi mức Importance cố định, Pruning trải trên một khoảng (thanh xám) thay vì một giá trị duy nhất.
  \item Do đó xác suất một Gaussian bị pruning tăng dần khi importance của nó giảm ($\sim 1-\text{importance}$ về xu hướng), nhưng đây là quan hệ \textbf{mang tính xác suất/thống kê}, không phải phép cắt tất định 1--1.
\end{itemize}
\centering
\includegraphics[width=0.78\linewidth]{3_6_pruning-vs-1-minus-imp.png}
\captionof{figure}{\footnotesize Trái: cùng Importance nhưng raw khác nhau tuỳ phân bố lỗi qua view; phải: khoảng giá trị Pruning ứng với mỗi mức Importance trên 300 Gaussian mô phỏng.}
\end{frame}
